% Options for packages loaded elsewhere
\PassOptionsToPackage{unicode}{hyperref}
\PassOptionsToPackage{hyphens}{url}
%
\documentclass[
  11pt,
]{article}
\usepackage{amsmath,amssymb}
\usepackage{iftex}
\ifPDFTeX
  \usepackage[T1]{fontenc}
  \usepackage[utf8]{inputenc}
  \usepackage{textcomp} % provide euro and other symbols
\else % if luatex or xetex
  \usepackage{unicode-math} % this also loads fontspec
  \defaultfontfeatures{Scale=MatchLowercase}
  \defaultfontfeatures[\rmfamily]{Ligatures=TeX,Scale=1}
\fi
\usepackage{lmodern}
\ifPDFTeX\else
  % xetex/luatex font selection
\fi
% Use upquote if available, for straight quotes in verbatim environments
\IfFileExists{upquote.sty}{\usepackage{upquote}}{}
\IfFileExists{microtype.sty}{% use microtype if available
  \usepackage[]{microtype}
  \UseMicrotypeSet[protrusion]{basicmath} % disable protrusion for tt fonts
}{}
\makeatletter
\@ifundefined{KOMAClassName}{% if non-KOMA class
  \IfFileExists{parskip.sty}{%
    \usepackage{parskip}
  }{% else
    \setlength{\parindent}{0pt}
    \setlength{\parskip}{6pt plus 2pt minus 1pt}}
}{% if KOMA class
  \KOMAoptions{parskip=half}}
\makeatother
\usepackage{xcolor}
\usepackage[margin=1in]{geometry}
\usepackage{longtable,booktabs,array}
\usepackage{calc} % for calculating minipage widths
% Correct order of tables after \paragraph or \subparagraph
\usepackage{etoolbox}
\makeatletter
\patchcmd\longtable{\par}{\if@noskipsec\mbox{}\fi\par}{}{}
\makeatother
% Allow footnotes in longtable head/foot
\IfFileExists{footnotehyper.sty}{\usepackage{footnotehyper}}{\usepackage{footnote}}
\makesavenoteenv{longtable}
\usepackage{graphicx}
\makeatletter
\def\maxwidth{\ifdim\Gin@nat@width>\linewidth\linewidth\else\Gin@nat@width\fi}
\def\maxheight{\ifdim\Gin@nat@height>\textheight\textheight\else\Gin@nat@height\fi}
\makeatother
% Scale images if necessary, so that they will not overflow the page
% margins by default, and it is still possible to overwrite the defaults
% using explicit options in \includegraphics[width, height, ...]{}
\setkeys{Gin}{width=\maxwidth,height=\maxheight,keepaspectratio}
% Set default figure placement to htbp
\makeatletter
\def\fps@figure{htbp}
\makeatother
\ifLuaTeX
  \usepackage{luacolor}
  \usepackage[soul]{lua-ul}
\else
  \usepackage{soul}
\fi
\setlength{\emergencystretch}{3em} % prevent overfull lines
\providecommand{\tightlist}{%
  \setlength{\itemsep}{0pt}\setlength{\parskip}{0pt}}
\setcounter{secnumdepth}{-\maxdimen} % remove section numbering
% Make all images max 0.5\textwidth
\usepackage{graphicx}
\setkeys{Gin}{width=0.5\textwidth,keepaspectratio}

% Modern fonts: Fira Sans headings, Fira Mono code
% Note: fontspec is already loaded by Pandoc, so we just set fonts
\setsansfont{FiraSans}[
    Path=/usr/share/texlive/texmf-dist/fonts/opentype/public/fira/,
    Extension=.otf,
    UprightFont=*-Regular,
    BoldFont=*-Bold,
    ItalicFont=*-Italic,
    BoldItalicFont=*-BoldItalic
]
\setmonofont{FiraMono}[
    Path=/usr/share/texlive/texmf-dist/fonts/opentype/public/fira/,
    Extension=.otf,
    Scale=0.90,
    UprightFont=*-Regular,
    BoldFont=*-Bold
]

% Make headings much larger with more spacing and sans-serif
\usepackage{titlesec}
\newcommand{\sectionbreak}{\clearpage}
\titleformat{\section}{\sffamily\LARGE\bfseries}{\thesection}{1em}{}
\titlespacing*{\section}{0pt}{24pt}{12pt}
\titleformat{\subsection}{\sffamily\Large\bfseries}{\thesubsection}{1em}{}
\titlespacing*{\subsection}{0pt}{20pt}{10pt}
\titleformat{\subsubsection}{\sffamily\large\bfseries}{\thesubsubsection}{1em}{}
\titlespacing*{\subsubsection}{0pt}{16pt}{8pt}

% Make document title Fira Sans Bold and left-aligned
\usepackage{titling}
\pretitle{\begin{flushleft}\sffamily\bfseries\huge}
\posttitle{\end{flushleft}}
\preauthor{\begin{flushleft}}
\postauthor{\end{flushleft}}
\predate{\begin{flushleft}}
\postdate{\end{flushleft}}

% Page header with footer content
\usepackage{fancyhdr}
\pagestyle{fancy}
\fancyhf{}
\fancyhead[L]{\small Elektrotechnik – Straub}
\fancyhead[R]{\small\itshape\nouppercase{\leftmark}}
\fancyfoot[C]{\small\thepage}
\renewcommand{\headrulewidth}{0.4pt}
\renewcommand{\footrulewidth}{0pt}
% Empty style for title page
\fancypagestyle{plain}{
  \fancyhf{}
  \renewcommand{\headrulewidth}{0pt}
  \renewcommand{\footrulewidth}{0pt}
}

% Increase paragraph spacing and line spacing
\usepackage{setspace}
\setstretch{1.3}
\setlength{\parskip}{0.8em}

% Use ragged right (left-aligned) instead of justified text
\raggedright

% More spacing around lists
\usepackage{enumitem}
\setlist{itemsep=0.5em,parsep=0.3em}

% Add moderate padding to code blocks
\usepackage{xcolor}
\definecolor{shadecolor}{RGB}{248,248,248}
\usepackage{fancyvrb}
\usepackage{framed}
\RecustomVerbatimEnvironment{Verbatim}{Verbatim}{fontsize=\small}
\ifLuaTeX
  \usepackage{selnolig}  % disable illegal ligatures
\fi
\usepackage{bookmark}
\IfFileExists{xurl.sty}{\usepackage{xurl}}{} % add URL line breaks if available
\urlstyle{same}
\hypersetup{
  pdftitle={Elektrotechnik (Teil 1/3)},
  hidelinks,
  pdfcreator={LaTeX via pandoc}}

\title{Elektrotechnik (Teil 1/3)}
\author{}
\date{}

\begin{document}
\maketitle

\textbf{Luft- und Raumfahrttechnik Bachelor, 1. Semester}

David Straub

\subsection{Organisatorisches}\label{organisatorisches}

\begin{itemize}
\tightlist
\item
  🎓 Moodle-Kurs: https://link.hm.edu/hnyu
\item
  💬 Matrix-Raum: https://link.hm.edu/83kf
\item
  🕥 Sprechstunde: Do. 10:30--11:30, Raum, B 374
\item
  📖 Literatur

  \begin{itemize}
  \tightlist
  \item
    Pregla -- \href{https://link.hm.edu/2c6h}{OPAC}
  \item
    Hagmann -- \href{https://link.hm.edu/fvqd}{OPAC}
  \item
    Hering u.a. --
    \href{https://link.springer.com/book/10.1007/978-3-662-67538-0}{online}
  \item
    Fischer --
    \href{https://link.springer.com/book/10.1007/978-3-658-25644-9}{online}
  \end{itemize}
\item
  🗒️ Vorlesungsskript Prof.~Palme u.a.: https://palme.userweb.mwn.de/
\item
  ✍️ Prüfung: schriftlich, 60 Minuten, keine Hilfsmittel
\end{itemize}

\subsection{Gliederung des Kurses}\label{gliederung-des-kurses}

\begin{enumerate}
\def\labelenumi{\arabic{enumi}.}
\tightlist
\item
  \hyperref[einfuxfchrung]{Einführung} (Physikalische Größen, Einheiten)
\item
  \hyperref[das-elektrische-feld]{Das elektrische Feld} (Ladungen,
  Kräfte, Felder, Potential, Spannung, Kapazität, Kondensatoren)
\item
  \hyperref[gleichstrom]{Gleichstrom} (Stromstärke, Widerstand,
  Stromkreisberechnungen, Energie, Leistung)
\item
  \hyperref[magnetismus]{Magnetismus} (Feld in Vakuum und Materie,
  Kräfte, magnetischer Kreis)
\item
  \hyperref[elektromagnetische-induktion]{Elektromagnetische Induktion}
  (Induktion, Selbstinduktion, Energie)
\item
  \hyperref[wechselstrom]{Wechselstrom} (Komplexe Wechselstromrechnung,
  Schaltungen, Leistung)
\item
  \hyperref[drehstrom]{Drehstrom} (Dreiphasensystem)
\item
  \hyperref[schaltvorguxe4nge-an-kapazituxe4ten-und-induktivituxe4ten]{Schaltvorgänge
  an Kapazitäten und Induktivitäten}
\end{enumerate}

\section{Einführung}\label{einfuxfchrung}

\begin{enumerate}
\def\labelenumi{\arabic{enumi}.}
\tightlist
\item
  Physikalische Größen
\item
  Internationales Einheitensystem (SI)
\end{enumerate}

\subsection{Physikalische Größen}\label{physikalische-gruxf6uxdfen}

\ldots{} sind messbare Eigenschaften eines Systems.

\textbf{Skalare Größen}: werden durch einen \emph{Zahlenwert} und eine
\emph{Einheit} beschrieben.

\[x = \underbrace{\lbrace x \rbrace}_{\text{Zahlenwert}} \cdot \underbrace{[x]}_{\text{Einheit}}\]

Beispiele:

\begin{itemize}
\tightlist
\item
  \(t = 10 \, \text{s}\) (Zeit)
\item
  \(m = 5 \, \text{kg}\) (Masse)
\item
  \(\Delta T = -20 \, \text{K}\) (Temperaturdifferenz)
\end{itemize}

\subsection{Rechnen mit Einheiten}\label{rechnen-mit-einheiten}

\begin{itemize}
\tightlist
\item
  Nur Größen mit gleichen Einheiten können addiert oder subtrahiert
  werden
\end{itemize}

\[x = 2 \, \text{m} + 3 \, \text{m} = 5 \, \text{m}\]

\begin{itemize}
\tightlist
\item
  Bei Multiplikation/Division von Größen werden die Einheiten
  multipliziert/dividiert
\end{itemize}

\[v = \frac{s}{t} = \frac{10 \, \text{m}}{5 \, \text{s}} = 2 \, \frac{\text{m}}{\text{s}} = 7{,}2 \, \text{km/h}\]

Hinweis: im Textsatz werden Einheiten immer aufrecht geschrieben,
Variablen \emph{kursiv}.

\subsection{Vektorielle physikalische
Größen}\label{vektorielle-physikalische-gruxf6uxdfen}

\ldots{} sind physikalische Größen, die durch einen \emph{Betrag} und
eine \emph{Richtung} beschrieben werden. Der Betrag wird durch einen
\emph{Zahlenwert} und eine \emph{Einheit} beschrieben.

\[\vec{v}\equiv \mathbf{v} = \underbrace{|\vec{v}|}_{\text{Betrag}} \cdot \underbrace{\vec{e}_v}_{\text{Richtung}}\]

\[ |\vec{v}| \equiv v = \underbrace{\lbrace v \rbrace}_{\text{Zahlenwert}} \cdot \underbrace{[v]}_{\text{Einheit}}\]

Der Zahlenwert des Betrags ist immer positiv.

Beispiele:

\begin{itemize}
\tightlist
\item
  \(\vec{v} = 10 \, \frac{\text{m}}{\text{s}} \cdot \vec{e}_x\)
  (Geschwindigkeit)
\item
  \(\vec{a} = 9{,}81 \, \frac{\text{m}}{\text{s}^2} \cdot (\vec{e}_{-z})\)
  (Beschleunigung)
\end{itemize}

\subsection{Das Internationale Einheitensystem
(SI)}\label{das-internationale-einheitensystem-si}

\begin{longtable}[]{@{}
  >{\raggedright\arraybackslash}p{(\columnwidth - 8\tabcolsep) * \real{0.3053}}
  >{\raggedright\arraybackslash}p{(\columnwidth - 8\tabcolsep) * \real{0.1789}}
  >{\raggedright\arraybackslash}p{(\columnwidth - 8\tabcolsep) * \real{0.2526}}
  >{\raggedright\arraybackslash}p{(\columnwidth - 8\tabcolsep) * \real{0.0947}}
  >{\raggedright\arraybackslash}p{(\columnwidth - 8\tabcolsep) * \real{0.1684}}@{}}
\toprule\noalign{}
\begin{minipage}[b]{\linewidth}\raggedright
Basisgröße
\end{minipage} & \begin{minipage}[b]{\linewidth}\raggedright
Größensymbol
\end{minipage} & \begin{minipage}[b]{\linewidth}\raggedright
Dimensionssymbol
\end{minipage} & \begin{minipage}[b]{\linewidth}\raggedright
Einheit
\end{minipage} & \begin{minipage}[b]{\linewidth}\raggedright
Einheitenzeichen
\end{minipage} \\
\midrule\noalign{}
\endhead
\bottomrule\noalign{}
\endlastfoot
Zeit & \(t\) & \(\text{T}\) & Sekunde & s \\
Länge & \(l\) & \(\text{L}\) & Meter & m \\
Masse & \(m\) & \(\text{M}\) & Kilogramm & kg \\
Elektrische Stromstärke & \(I\) & \(\text{I}\) & Ampere & A \\
Thermodynamische Temperatur & \(T\) & \(\Theta\) & Kelvin & K \\
Stoffmenge & \(n\) & \(\text{N}\) & Mol & mol \\
Lichtstärke & \(I_v\) & \(\text{J}\) & Candela & cd \\
\end{longtable}

\subsection{Naturkonstanten und
SI-Einheiten}\label{naturkonstanten-und-si-einheiten}

\begin{longtable}[]{@{}
  >{\raggedright\arraybackslash}p{(\columnwidth - 6\tabcolsep) * \real{0.1414}}
  >{\raggedright\arraybackslash}p{(\columnwidth - 6\tabcolsep) * \real{0.5455}}
  >{\raggedright\arraybackslash}p{(\columnwidth - 6\tabcolsep) * \real{0.2222}}
  >{\raggedright\arraybackslash}p{(\columnwidth - 6\tabcolsep) * \real{0.0909}}@{}}
\toprule\noalign{}
\begin{minipage}[b]{\linewidth}\raggedright
Konstante
\end{minipage} & \begin{minipage}[b]{\linewidth}\raggedright
Beschreibung
\end{minipage} & \begin{minipage}[b]{\linewidth}\raggedright
Exakter Wert
\end{minipage} & \begin{minipage}[b]{\linewidth}\raggedright
Einheit
\end{minipage} \\
\midrule\noalign{}
\endhead
\bottomrule\noalign{}
\endlastfoot
\(\Delta\nu_\mathrm{Cs}\) & Strahlung des Caesium-Atoms &
9\,192\,631\,770 & Hz \\
\(c\) & Lichtgeschwindigkeit & 299\,792\,458 & m/s \\
\(h\) & Planck-Konstante & 6,62607015\,×\,10−34 & J·s \\
\(e\) & Elementarladung & 1,602176634\,×\,10−19 & C \\
\(k_\mathrm{B}\) & Boltzmann-Konstante & 1,380649\,×\,10−23 & J/K \\
\(N_\mathrm{A}\) & Avogadro-Konstante & 6,02214076\,×\,1023 & mol⁻¹ \\
\(K_\mathrm{cd}\) & Photometrisches Strahlungsäquivalent & 683 & lm/W \\
\end{longtable}

\subsection{Abgeleitete Einheiten}\label{abgeleitete-einheiten}

Von den Basisgrößen lassen sich durch mathematische Operationen
abgeleitete Einheiten bilden. Beispiele für abgeleitete Einheiten:

\begin{itemize}
\item
  \textbf{Kraft}: \(\vec{F} = m \cdot \vec{a}\)
  \([F] = [m] \cdot [\vec{a}]= \text{kg} \cdot \frac{\text{m}}{\text{s}^2} = \text{N}\)
  (Newton)
\item
  \textbf{Energie/Arbeit}: \(W = F \cdot s\)
  \([W]  = \text{N} \cdot \text{m} = \frac{\text{kg} \cdot \text{m}^2}{\text{s}^2}= \text{J}\)
  (Joule)
\item
  \textbf{Leistung}: \(P = \frac{\Delta W}{\Delta t}\)
  \([P]  = \frac{[W]}{[t]} = \frac{\text{J}}{\text{s}} = \frac{\text{kg} \cdot \text{m}^2}{\text{s}^3}= \text{W}\)
  (Watt)
\end{itemize}

\subsection{Dimensionsanalyse}\label{dimensionsanalyse}

Jede physikalische Größe hat -- unabängig von Einheit oder Zahlenwert --
eine \textbf{Dimension}, die beschreibt, wie die Größe aus den
Grundgrößen zusammengesetzt ist.

Beispiele:

\begin{itemize}
\tightlist
\item
  Geschwindigkeit: \(\text{dim}[v] = \frac{\text{L}}{\text{T}}\)
\item
  Kraft: \(\text{dim}[F] = \text{M} \cdot \frac{\text{L}}{\text{T}^2}\)
\item
  Winkel: \(\text{dim}[\varphi] = \frac{\text{L}}{\text{L}} = 1\)
  (dimensionslos)
\end{itemize}

Beide Seiten einer Gleichung müssen dieselbe Dimension haben!

\subsection{SI-Präfixe (alltäglich)}\label{si-pruxe4fixe-alltuxe4glich}

\begin{longtable}[]{@{}
  >{\raggedright\arraybackslash}p{(\columnwidth - 10\tabcolsep) * \real{0.1571}}
  >{\raggedright\arraybackslash}p{(\columnwidth - 10\tabcolsep) * \real{0.1143}}
  >{\raggedright\arraybackslash}p{(\columnwidth - 10\tabcolsep) * \real{0.2286}}
  >{\raggedright\arraybackslash}p{(\columnwidth - 10\tabcolsep) * \real{0.1571}}
  >{\raggedright\arraybackslash}p{(\columnwidth - 10\tabcolsep) * \real{0.1143}}
  >{\raggedright\arraybackslash}p{(\columnwidth - 10\tabcolsep) * \real{0.2286}}@{}}
\toprule\noalign{}
\begin{minipage}[b]{\linewidth}\raggedright
Faktor
\end{minipage} & \begin{minipage}[b]{\linewidth}\raggedright
Name
\end{minipage} & \begin{minipage}[b]{\linewidth}\raggedright
Präfix
\end{minipage} & \begin{minipage}[b]{\linewidth}\raggedright
Faktor
\end{minipage} & \begin{minipage}[b]{\linewidth}\raggedright
Name
\end{minipage} & \begin{minipage}[b]{\linewidth}\raggedright
Präfix
\end{minipage} \\
\midrule\noalign{}
\endhead
\bottomrule\noalign{}
\endlastfoot
\(10^{-1}\) & Dezi & d & \(10^{1}\) & Deka & da \\
\(10^{-2}\) & Zenti & c & \(10^{2}\) & Hekto & h \\
\(10^{-3}\) & Milli & m & \(10^{3}\) & Kilo & k \\
\(10^{-6}\) & Mikro & µ & \(10^{6}\) & Mega & M \\
\(10^{-9}\) & Nano & n & \(10^{9}\) & Giga & G \\
\(10^{-12}\) & Piko & p & \(10^{12}\) & Tera & T \\
\end{longtable}

\subsection{SI-Präfixe (nicht
alltäglich)}\label{si-pruxe4fixe-nicht-alltuxe4glich}

\begin{longtable}[]{@{}
  >{\raggedright\arraybackslash}p{(\columnwidth - 10\tabcolsep) * \real{0.1571}}
  >{\raggedright\arraybackslash}p{(\columnwidth - 10\tabcolsep) * \real{0.1143}}
  >{\raggedright\arraybackslash}p{(\columnwidth - 10\tabcolsep) * \real{0.2286}}
  >{\raggedright\arraybackslash}p{(\columnwidth - 10\tabcolsep) * \real{0.1571}}
  >{\raggedright\arraybackslash}p{(\columnwidth - 10\tabcolsep) * \real{0.1143}}
  >{\raggedright\arraybackslash}p{(\columnwidth - 10\tabcolsep) * \real{0.2286}}@{}}
\toprule\noalign{}
\begin{minipage}[b]{\linewidth}\raggedright
Faktor
\end{minipage} & \begin{minipage}[b]{\linewidth}\raggedright
Name
\end{minipage} & \begin{minipage}[b]{\linewidth}\raggedright
Präfix
\end{minipage} & \begin{minipage}[b]{\linewidth}\raggedright
Faktor
\end{minipage} & \begin{minipage}[b]{\linewidth}\raggedright
Name
\end{minipage} & \begin{minipage}[b]{\linewidth}\raggedright
Präfix
\end{minipage} \\
\midrule\noalign{}
\endhead
\bottomrule\noalign{}
\endlastfoot
\(10^{-15}\) & Femto & f & \(10^{15}\) & Peta & P \\
\(10^{-18}\) & Atto & a & \(10^{18}\) & Exa & E \\
\(10^{-21}\) & Zepto & z & \(10^{21}\) & Zetta & Z \\
\(10^{-24}\) & Yokto & y & \(10^{24}\) & Yotta & Y \\
\(10^{-27}\) & Ronto & r & \(10^{27}\) & Ronna & R \\
\(10^{-30}\) & Quecto & q & \(10^{30}\) & Quetta & Q \\
\end{longtable}

\subsection{µ \& °C: praktische Tipps}\label{uxb5-c-praktische-tipps}

\begin{itemize}
\tightlist
\item
  Mikro: µ (griechischer Buchstabe ``My'')

  \begin{itemize}
  \tightlist
  \item
    Deutsches Tastaturlayout: \texttt{AltGr} + \texttt{m}
  \end{itemize}
\item
  Grad Celsius: °C (Gradzeichen + Großbuchstabe C)

  \begin{itemize}
  \tightlist
  \item
    Deutsches Tastaturlayout: \texttt{Shift} + \texttt{\^{}} +
    \texttt{C}
  \end{itemize}
\end{itemize}

Nur in Systemen, die diese Schriftzeichen nicht unterstützen (ASCII)
laut DIN 66030:2002-05:

\begin{itemize}
\tightlist
\item
  µ → u
\item
  °C → Cel
\end{itemize}

\subsection{⚠️ Nicht-SI-Einheiten in der Luftfahrt
✈️}\label{nicht-si-einheiten-in-der-luftfahrt}

Immer noch weit verbreitet:

\begin{itemize}
\tightlist
\item
  Flughöhe in \textbf{Fuß} 🦶

  \begin{itemize}
  \tightlist
  \item
    1 ft = 0,3048 m
  \end{itemize}
\item
  Entfernung in \textbf{Seemeilen} 🚢

  \begin{itemize}
  \tightlist
  \item
    1 NM = 1852 m
  \end{itemize}
\item
  Geschwindigkeit in \textbf{Knoten} 🪢

  \begin{itemize}
  \tightlist
  \item
    1 kt = 1 NM/h = 1,852 km/h
  \end{itemize}
\end{itemize}

\includegraphics{/tmp/tmp.sFQRuqDRHE/images/image_1.pdf}

\subsection{Mars Climate Orbiter}\label{mars-climate-orbiter}

https://www.youtube.com/watch?v=MfavzjbZzl8

\section{Das elektrische Feld}\label{das-elektrische-feld}

\begin{enumerate}
\def\labelenumi{\arabic{enumi}.}
\tightlist
\item
  Elektrische Ladung
\item
  Coulomb'sches Gesetz
\item
  Elektrisches Feld im Vakuum
\item
  \hyperref[feldlinien]{Feldlinien und Gauß'sches Gesetz}
\item
  Elektrisches Feld in Materie
\item
  Potential, Spannung, Arbeit
\item
  Homogenes Feld und Kondensatoren
\end{enumerate}

\subsection{Die vier fundamentalen
Wechselwirkungen}\label{die-vier-fundamentalen-wechselwirkungen}

\begin{enumerate}
\def\labelenumi{\arabic{enumi}.}
\tightlist
\item
  \textbf{Gravitation} 🪐

  \begin{itemize}
  \tightlist
  \item
    Hält das Sonnensystem zusammen -- wirkt auf Masse
  \end{itemize}
\item
  \textbf{Elektromagnetismus} ⚡

  \begin{itemize}
  \tightlist
  \item
    Hält Atome und Moleküle zusammen -- wirkt auf elektrische Ladung
  \end{itemize}
\item
  \textbf{Starke Wechselwirkung} 🎨

  \begin{itemize}
  \tightlist
  \item
    Hält Atomkerne zusammen
  \end{itemize}
\item
  \textbf{Schwache Wechselwirkung} ☢️

  \begin{itemize}
  \tightlist
  \item
    Verantwortlich für radioaktiven Zerfall
  \end{itemize}
\end{enumerate}

\subsection{}\label{section}

\includegraphics{/tmp/tmp.sFQRuqDRHE/images/dc85f6a65900a13d052e17a62908ce6f.jpg}
\includegraphics{/tmp/tmp.sFQRuqDRHE/images/3748ce0af9bdae4242372678673a81d1.jpg}

\subsection{\texorpdfstring{Elektrische Ladung (\emph{electric
charge})}{Elektrische Ladung (electric charge)}}\label{elektrische-ladung-electric-charge}

\begin{itemize}
\tightlist
\item
  Alle Materie besteht aus Elementarteilchen, von denen einige
  elektrische Ladungen tragen
\item
  Elektrische Ladungen treten in zwei Arten auf: positive und negative
  Ladungen (Vorzeichen: Konvention!)
\item
  Gleichnamige Ladungen stoßen sich ab, ungleichnamige ziehen sich an
\end{itemize}

\subsection{Aufbau der Materie}\label{aufbau-der-materie}

\begin{itemize}
\tightlist
\item
  Atome bestehen aus positiv geladenen Protonen, neutralen Neutronen und
  negativ geladenen Elektronen
\item
  Protonen und Neutronen bilden den Atomkern
\item
  Elektronen bewegen sich in der Atomhülle um den Atomkern
\end{itemize}

\includegraphics{/tmp/tmp.sFQRuqDRHE/images/de142681894db369bf949e141c52f92f.png}

\subsection{Elementarladung}\label{elementarladung}

\begin{itemize}
\tightlist
\item
  Elektrische Ladungen sind immer ganzzahlige Vielfache der
  Elementarladung \(e= 1{,}602176634 \cdot 10^{-19} \, \text{C}\)
  (Definition des Coulombs)

  \begin{itemize}
  \tightlist
  \item
    Elektron: \(Q = -e\) (!)
  \item
    Proton: \(Q = +e\)
  \item
    Up-Quark: \(Q = +\frac{2}{3}e\), Down-Quark: \(Q = -\frac{1}{3}e\)
  \end{itemize}
\item
  Man sagt, die Ladung sei \emph{quantisiert}
\end{itemize}

\includegraphics{/tmp/tmp.sFQRuqDRHE/images/image_5.pdf}

\subsection{Coulumb'sches Gesetz --
Experiment}\label{coulumbsches-gesetz-experiment}

https://www.youtube.com/watch?v=9mFlELwuctI

\subsection{\texorpdfstring{Coulomb'sches Gesetz (\emph{Coulomb's
law})}{Coulomb'sches Gesetz (Coulomb's law)}}\label{coulombsches-gesetz-coulombs-law}

\begin{itemize}
\tightlist
\item
  Experimente haben gezeigt, dass die Kraft zwischen zwei Punktladungen
  \(Q_1\) und \(Q_2\) proportional zur Größe der Ladungen und umgekehrt
  proportional zum Quadrat ihres Abstands \(r\) ist
\item
  Mathematisch wird dies durch das Coulomb'sche Gesetz beschrieben:
\end{itemize}

\[|\vec{F}_{12}| = k \cdot \frac{Q_1 \cdot Q_2}{r^2}\]

\(k\) ist abhängig vom Einheitensystem. Im SI-System gilt:
\(k = \frac{1}{4 \pi \varepsilon_0}\), wobei \(\varepsilon_0\) die
elektrische Feldkonstante ist,
\(\varepsilon_0 \approx 8{,}854 \cdot 10^{-12} \, \frac{\text{C}^2}{\text{N} \cdot \text{m}^2}\).

\subsection{Analogie zur Schwerkraft}\label{analogie-zur-schwerkraft}

Newtonsches Gravitationsgesetz: Kraft zwischen zwei Himmelskörpern

\[|\vec{F}_{12}| = G \cdot \frac{m_1 \cdot m_2}{r^2}\]

\(G\): Gravitationskonstante,
\(G \approx 6{,}6743 \cdot 10^{-11} \, \frac{\text{kg}^2}{\text{N} \cdot \text{m}^2}\)

\subsection{Beispiel: Relative Stärke von Coulomb- und
Gravitationskraft}\label{beispiel-relative-stuxe4rke-von-coulomb--und-gravitationskraft}

\begin{itemize}
\tightlist
\item
  Proton: \(m_p \approx 1{,}67 \cdot 10^{-27} \, \text{kg}\),
  \(Q_p = +e\)
\item
  Elektron: \(m_e \approx 9{,}11 \cdot 10^{-31} \, \text{kg}\),
  \(Q_e = -e\)
\item
  \(|\vec{F}_{12,\text{Coulomb}}| = \frac{1}{4 \pi \varepsilon_0} \cdot \frac{Q_1 \cdot Q_2}{r^2}\)
\item
  \(\epsilon_0 \approx 8{,}854 \cdot 10^{-12} \, \frac{\text{As}}{\text{Vm}}\)
\item
  \(|\vec{F}_{12,\text{Gravitation}}| = G \cdot \frac{m_1 \cdot m_2}{r^2}\)
\item
  \(G \approx 6{,}6743 \cdot 10^{-11} \, \frac{\text{m}^3}{\text{kg} \cdot \text{s}^2}\)
\end{itemize}

\subsection{Elektromagnetismus im
Alltag}\label{elektromagnetismus-im-alltag}

Fast alle alltäglichen physikalischen Phänomene werden von der
elektromagnetischen Wechselwirkung bestimmt!

Die Gravitation spielt nur eine Rolle, da

\begin{itemize}
\tightlist
\item
  es keine negativen Massen gibt → immer anziehend
\item
  die elektrischen Ladungen von Elektronen und Protonen exakt aufheben
\end{itemize}

\subsection{\texorpdfstring{Elektrische Feldstärke (\emph{electric field
{[}strength{]}})}{Elektrische Feldstärke (electric field {[}strength{]})}}\label{elektrische-feldstuxe4rke-electric-field-strength}

\begin{itemize}
\tightlist
\item
  Ein elektrisch geladenes Teilchen übt eine Kraft auf andere elektrisch
  geladene Teilchen aus
\item
  Diese Kraft ist umso größer, je größer die Ladung der Probeteilchen
  ist
\item
  Elektrische Feldstärke: Kraft pro Ladungseinheit, die auf eine
  Probeladung wirkt
\end{itemize}

\[\vec{E} = \frac{\vec{F}}{Q} \Leftrightarrow \vec{F} = Q \cdot \vec{E}\]

Feld = ortsabhängige physikalische Größe (Vektorfeld/Skalarfeld)

\([\vec{E}] = \frac{\text{N}}{\text{C}}\)

\subsection{Elektrisches Feld einer
Punktladung}\label{elektrisches-feld-einer-punktladung}

Die elektrische Feldstärke \(\vec{E}\) im Abstand \(r=|\vec{r}|\) einer
Punktladung \(Q\) ist:

\[\vec{E}(\vec r) = \frac{Q}{4 \cdot \pi \cdot \varepsilon_0 \cdot r^2} \cdot \frac{\vec{r}}{r} = \frac{Q}{4 \cdot \pi \cdot \varepsilon_0 \cdot r^2} \cdot \vec{e}_r\]

\subsection{Feldlinien}\label{feldlinien}

\begin{itemize}
\tightlist
\item
  Die Feldlinien eines elektrischen Feldes sind Linien, die die Richtung
  und Stärke des Feldes darstellen
\item
  Sie verlaufen von positiven zu negativen Ladungen und zeigen die
  Richtung der Kraft an, die auf eine positive Probeladung wirken würde
\item
  Die Dichte der Feldlinien ist proportional zur Stärke des elektrischen
  Feldes: Je dichter die Linien, desto stärker das Feld
\end{itemize}

\subsection{Superpositionsprinzip}\label{superpositionsprinzip}

Das elektrische Feld mehrerer (diskreter) Ladungen ist die Vektorsumme
der Felder der einzelnen Ladungen

\[\vec{E}(\vec{r}) = \sum_{i} \vec{E}_i(\vec{r})\]
\[= \sum_{i} \frac{Q_i}{4 \pi \varepsilon_0 \cdot |\vec{r}-\vec{r}_i|^2} \cdot \frac{\vec{r}-\vec{r}_i}{|\vec{r}-\vec{r}_i|}\]

Übergang zu kontinuierlicher Ladungsverteilung: Integral

\[\vec{E}(\vec{r}) = \int_{V} \frac{\rho(\vec{r}') \cdot dV'}{4 \pi \varepsilon_0 \cdot |\vec{r}-\vec{r}'|^2} \cdot \frac{\vec{r}-\vec{r}'}{|\vec{r}-\vec{r}'|}\]

\subsection{Feldlinien: Beispiele}\label{feldlinien-beispiele}

\begin{itemize}
\tightlist
\item
  Punktladung
\item
  Zwei gegensätzliche Punktladungen (Dipol)
\item
  Zwei gleichnamige Punktladungen
\item
  Positiv geladene Ebene
\item
  Zwei entgegengesetzt geladene Platten (Plattenkondensator)
\end{itemize}

\subsection{\texorpdfstring{Elektrische Flussdichte (\emph{electric flux
density})}{Elektrische Flussdichte (electric flux density)}}\label{elektrische-flussdichte-electric-flux-density}

\[\vec{D} = \varepsilon_0 \cdot \vec{E}\]

\[[\vec{D}] = [\varepsilon_0] \cdot [\vec{E}] = \frac{\text{C}^2}{\text{N} \cdot \text{m}^2} \cdot \frac{\text{N}}{\text{C}} = \frac{\text{C}}{\text{m}^2}\]

Fluss durch eine Fläche \(A\):

\[\Psi = \int_{A} \vec{D} \cdot d\vec{A}\]

\includegraphics{/tmp/tmp.sFQRuqDRHE/images/85cf0e3499426285a82cc26a94274622.jpg}

\subsection{Elektrische Flussdichte einer
Punktladung}\label{elektrische-flussdichte-einer-punktladung}

Die elektrische Flussdichte \(\vec{D}\) im Abstand \(r\) einer
Punktladung \(Q\) ist:

\[\vec{D}(r) = \frac{Q}{4 \cdot \pi \cdot r^2} \cdot \frac{\vec{r}}{r} = \frac{Q}{A(r)} \cdot \vec{e}_r\]

Für eine konstante Flussdichte \(D\) auf der Fläche \(A\) gilt:
\(Q = D(r) \cdot A(r)\)

\subsection{Fluss durch geschlossene
Oberflächen}\label{fluss-durch-geschlossene-oberfluxe4chen}

\subsection{\texorpdfstring{Satz von Gauß (\emph{Gauss's
law})}{Satz von Gauß (Gauss's law)}}\label{satz-von-gauuxdf-gausss-law}

Der elektrische Fluss durch eine geschlossene Oberfläche ist gleich der
eingeschlossenen Ladung:

\[\oint_{A} \vec{D} \cdot d\vec{A} = Q_{\text{innen}}\]

\subsection{Beispiel: Kugelsymmetrische
Ladungsverteilung}\label{beispiel-kugelsymmetrische-ladungsverteilung}

Für eine Punktladung \(Q\) im Zentrum einer Kugel mit Radius \(r\) gilt:

\[\oint_{A} \vec{D} \cdot d\vec{A} = D(r) \cdot 4\pi r^2 = Q\]

Daraus folgt:

\[D = \frac{Q}{4\pi r^2}\]

\subsection{Beispiel: Unendlich langer, gerader
Leiter}\label{beispiel-unendlich-langer-gerader-leiter}

Für einen unendlich langen Leiter mit Linienladungsdichte \(\lambda\)
verwenden wir eine zylindrische Oberfläche (Radius \(r\), Länge \(l\)).

\textbf{Symmetrie:} \(\vec{D}\) ist radial und konstant auf der
Mantelfläche

\[\oint_{A} \vec{D} \cdot d\vec{A} = D(r) \cdot 2\pi r l = \lambda l\]

\textbf{Ergebnis:}
\[D(r) = \frac{\lambda}{2\pi r}, \quad E(r) = \frac{\lambda}{2\pi \varepsilon_0 r}\]

\subsection{Elektrisches Feld in
Materie}\label{elektrisches-feld-in-materie}

\begin{itemize}
\tightlist
\item
  In nicht oder schwach leitenden Materialien können elektrische Felder
  zu \textbf{Polarisation} führen.
\item
  Bei der Polarisation richten sich die positiven und negativen Ladungen
  innerhalb des Materials unter dem Einfluss des elektrischen Feldes neu
  aus.
\item
  Dies führt zu einer Verschiebung der Ladungszentren und erzeugt ein
  internes elektrisches Feld, das dem äußeren Feld entgegenwirkt.
\item
  Solche polarisierbaren Materialien nennt man \textbf{Dielektrika}.
\end{itemize}

\subsection{Abschwächung des elektrischen Feldes in
Dielektrika}\label{abschwuxe4chung-des-elektrischen-feldes-in-dielektrika}

Definition:

\[\vec{E}_\text{Materie} = \underbrace{\vec{E}_\text{Vakuum}}_{\text{freie Ladungen}} - \underbrace{\vec{E}_\text{Polarisation}}_{\text{gebundene Ladungen}} = \frac{1}{\varepsilon_r} \cdot \vec{E}_\text{Vakuum}\]

mit \(\varepsilon_r \geq 1\) Permittivitätszahl (auch relative
Permittivität). Beispiele:

\begin{longtable}[]{@{}ll@{}}
\toprule\noalign{}
Material & \(\varepsilon_r\) \\
\midrule\noalign{}
\endhead
\bottomrule\noalign{}
\endlastfoot
Luft & 1,00059 \\
Gummi & 2,5--3,5 \\
Glas & 5--7 \\
Destilliertes Wasser & 81 \\
\end{longtable}

\subsection{Elektrische Flussdichte in
Dielektrika}\label{elektrische-flussdichte-in-dielektrika}

\textbf{Konvention}: man vereinbart, dass die elektrische Flussdichte
sich immer auf das durch die freien Ladungen erzeugte Feld bezieht.

\[\vec{D} =\varepsilon_0 \vec{E}_\text{Materie} + \vec{P} = \varepsilon_0 \varepsilon_r \vec{E}_\text{Materie} = 
\varepsilon \vec{E}_\text{Materie} = \varepsilon_0 \cdot \vec{E}_\text{Vakuum}\]

Polarisation \(\vec{P}= -\varepsilon_0 \vec{E}_\text{Polarisation}\)

\(\vec{D}\) heißt daher auch (veraltet) \textbf{Elektrische
Verschiebungsdichte} (\emph{electric displacement field}).

Vorteil: der Satz von Gauß gilt unverändert, wenn man nur die freien
Ladungen berücksichtigt:

\[\oint_{A} \vec{D} \cdot d\vec{A} = Q_{\text{innen, frei}}\]

\subsection{Elektrische Arbeit}\label{elektrische-arbeit}

Bewegung einer positiven Probeladung \(Q_P\) im Feld einer positiven
Punktladung \(Q^+\)

\(P_1 \rightarrow P_2\): \(W > 0\) : Arbeit wird freigesetzt
\(P_2 \rightarrow P_1\) : \(W < 0\) : Arbeit muss aufgebracht werden

Vgl. Mechanik: \(W = \vec{F} \cdot \vec{s}\)

Hier: \(\vec{F}\) ist abhängig von \(r\)!

\[W_{12} = \sum_{i} \Delta W_{i} = \sum_{i} F_{i} \cdot \Delta r\]

\subsection{Elektrische Arbeit und
Potential}\label{elektrische-arbeit-und-potential}

\[W_{12} = \sum_{i} \Delta W_{i} = \sum_{i} F_{i} \cdot \Delta r\]

Integral: \(\Delta r \rightarrow 0\)

\[W_{12} = \int_{r_1}^{r_2} F(r) \cdot dr = Q_P \cdot \int_{r_1}^{r_2} E(r) \cdot dr
= Q_P \cdot \int_{r_1}^{r_2} \frac{Q^+}{4 \cdot \pi \cdot \varepsilon \cdot r^2} dr\]
\[= Q_P \cdot \left[ \frac{Q^+}{4 \cdot \pi \cdot \varepsilon} \cdot \frac{1}{r_1} - \frac{Q^+}{4 \cdot \pi \cdot \varepsilon} \cdot \frac{1}{r_2} \right]\]
\[= Q_P \cdot (\varphi_1 - \varphi_2)\]

\subsection{Elektrisches Potential einer
Punktladung}\label{elektrisches-potential-einer-punktladung}

Das elektrisches Potential \(\varphi\) im Abstand \(r\) von einer
Punktladung \(Q\):

\[\varphi(r) = \frac{Q}{4 \cdot \pi \cdot \varepsilon_0 \cdot \varepsilon_r \cdot r}\]

Einheit: \([\varphi] = \frac{\text{J}}{\text{C}} = \text{V}\) (Volt)

Punkte gleichen Potentials bilden Äquipotentialflächen.

\subsection{Analogie:
Gravitationspotential}\label{analogie-gravitationspotential}

Potentielle Energie bzgl. Referenzhöhe:

\[E_\text{pot} = m \cdot g \cdot h = m \cdot \varphi_g(h)\]

\begin{itemize}
\tightlist
\item
  Höhenlinien sind Äquipotentiallinien
\item
  Die Kraft wirkt immer in Richtung des stärksten Gefälles (senkrecht zu
  den Äquipotentiallinien)
\item
  Die „Feldlinien`` sind nie in sich geschlossen
\end{itemize}

\subsection{Potentialfelder}\label{potentialfelder}

Elektrostatische Felder sind \textbf{Potentialfelder} oder auch
\textbf{wirbelfreie Felder}. Für sie gilt:

\begin{itemize}
\tightlist
\item
  Feldlinien beginnen und enden auf Ladungen („Quellen`` oder „Senken``)
\item
  Feldlinien sind nie in sich geschlossen
\item
  Das Feld lässt sich als Gradient (Richtungsableitung) eines
  Skalarfeldes (Potential) darstellen
\end{itemize}

Weitere Beispiele für Potentialfelder:

\begin{itemize}
\tightlist
\item
  Schwerkraft auf der Erdoberfläche (2D)
\item
  Schwerkraft zwischen Himmelskörpern (3D)
\item
  Wärmefluss in Festkörpern (1D, 2D, 3D)
\item
  Grundwasserstrom (3D) (nur solange die Strömung wirbelfrei ist!)
\end{itemize}

\subsection{Spannung \& Arbeit}\label{spannung-arbeit}

\subsubsection{\texorpdfstring{Elektrische Spannung (\emph{electric
voltage})}{Elektrische Spannung (electric voltage)}}\label{elektrische-spannung-electric-voltage}

Die elektrische Spannung ist definiert als Potentialdifferenz:

\[U_{12} = \varphi_1 - \varphi_2\]

Einheit: \([U] = \text{V}\) (Volt)

\subsubsection{\texorpdfstring{Elektrische Arbeit (\emph{electric field
work})}{Elektrische Arbeit (electric field work)}}\label{elektrische-arbeit-electric-field-work}

Die elektrische Arbeit ist das Produkt aus Ladung und Spannung:

\[W_{12} = \int_{P_1}^{P_2} \vec{F} \cdot d\vec{s} = Q \cdot (\varphi_1 - \varphi_2) = Q \cdot U_{12}\]

Einheit: \([W] = \text{J}\) (Joule)

Die elektrische Arbeit ist unabhängig vom Weg!

\subsection{Beziehung zwischen elektrischem Feld und
Spannung}\label{beziehung-zwischen-elektrischem-feld-und-spannung}

\[W_{12} = \int_{P_1}^{P_2} \vec{F} \cdot d\vec{s} = Q \cdot (\varphi_1 - \varphi_2) = Q \cdot U_{12}\]

Für die elektrische Spannung gilt allgemein:

\[U_{12} = \int_{P_1}^{P_2} \vec{E} \cdot d\vec{s} = \varphi_1 - \varphi_2\]

\subsection{Homogenes elektrisches
Feld}\label{homogenes-elektrisches-feld}

Ein homogenes elektrisches Feld ist durch konstante Feldstärke und
parallele Feldlinien gekennzeichnet.

Wichtige Eigenschaften:

\begin{itemize}
\tightlist
\item
  Konstante Feldstärke \(E\) in Betrag und Richtung
\item
  Parallele Feldlinien
\item
  Äquipotentialflächen stehen senkrecht zu den Feldlinien
\item
  Die Spannung zwischen zwei Punkten ist \(U = E \cdot d\), wobei \(d\)
  der Abstand in Feldrichtung ist
\end{itemize}

\includegraphics{/tmp/tmp.sFQRuqDRHE/images/a805089236c4c198a065d6cb9d1fe808.png}

\subsection{Homogenes Feld mit dem Satz von Gauß: linke
Seite}\label{homogenes-feld-mit-dem-satz-von-gauuxdf-linke-seite}

Unendlich ausgedehnte, gleichmäßig geladene Ebene mit
Flächenladungsdichte \(\sigma\)

\textbf{Gesucht:} Elektrische Feldstärke \(E\) im Abstand \(d\) von der
Ebene

\textbf{Ansatz:} Anwendung des Satzes von Gauß mit einem zylindrischen
Gauß'schen Volumen

\[\oint_{A} \vec{D} \cdot d\vec{A} = Q_{\text{innen}}\]

\textbf{Gauß'sche Fläche:} Zylinder mit Grundfläche \(A\) und Achse
senkrecht zur geladenen Ebene

\begin{itemize}
\tightlist
\item
  Mantelfläche:
  \(\vec{D} \perp d\vec{A} \Rightarrow \vec{D} \cdot d\vec{A} = 0\)
\item
  Grundflächen:
  \(\vec{D} \parallel d\vec{A} \Rightarrow \vec{D} \cdot d\vec{A} = D \cdot dA\)
\end{itemize}

\[\oint_{A} \vec{D} \cdot d\vec{A} = D \cdot A + D \cdot A = 2 \cdot D \cdot A\]

\subsection{Homogenes Feld mit dem Satz von Gauß: rechte
Seite}\label{homogenes-feld-mit-dem-satz-von-gauuxdf-rechte-seite}

\textbf{Eingeschlossene Ladung:} \(Q_{\text{innen}} = \sigma \cdot A\)

\textbf{Satz von Gauß:} \[2 \cdot D \cdot A = \sigma \cdot A\]

\[D = \frac{\sigma}{2} \quad \Rightarrow \quad E = \frac{\sigma}{2 \varepsilon_0 \varepsilon_r}\]

\textbf{Ergebnis:} Das Feld ist homogen und unabhängig vom Abstand zur
Ebene.

\subsection{\texorpdfstring{Kondensatoren
(\emph{capacitors})}{Kondensatoren (capacitors)}}\label{kondensatoren-capacitors}

Kondensatoren sind elektrische Bauelemente, die elektrische Ladung
speichern können.

Die gespeicherte Ladung für eine gegebene Spannung wird bezeichnet als:

\textbf{Kapazität (\emph{capacitance})}

\[C := \frac{Q}{U} \]

Einheit: \([C] = \frac{\text{C}}{\text{V}} = \text{F}\) (Farad)

\includegraphics{/tmp/tmp.sFQRuqDRHE/images/eda09bfc85a3947d780118d3285d197a.jpg}

\subsection{⚠️ Kapazität ≠ Kapazität}\label{kapazituxe4t-kapazituxe4t}

Die Kapazität (\emph{capacity}) einer Batterie ist eine Ladungsmenge!

z.B.:
\(\text{mAh} = 10^{-3} \, \text{A} \cdot 3600 \, \text{s} = 3{,}6 \, \text{C}\)

Nicht zu verwechseln mit der Kapazität (\emph{capacitance}) eines
Kondensators in Farad!

\subsection{Plattenkondensator}\label{plattenkondensator}

\[E =\frac{Q}{\varepsilon_0 \varepsilon_r A} ,\qquad U = E \cdot d = \frac{Q}{\varepsilon_0 \varepsilon_r A} \cdot d\]

\[C=\frac{Q}{U} = \frac{\varepsilon_0 \varepsilon_r A}{d}=\frac{\varepsilon A}{d}\]

Kapazität steigt mit:

\begin{itemize}
\tightlist
\item
  Fläche \(A\) der Platten
\item
  relativer Permittivität \(\varepsilon_r\) des Dielektrikums
\item
  Abnahme des Plattenabstands \(d\)
\end{itemize}

\includegraphics{/tmp/tmp.sFQRuqDRHE/images/image_9.pdf}

\subsection{Kugelkondensator}\label{kugelkondensator}

Ein Kugelkondensator besteht aus zwei konzentrischen leitenden
Kugelschalen mit den Radien \(R_i\) (innen) und \(R_a\) (außen).

\includegraphics{/tmp/tmp.sFQRuqDRHE/images/image_10.pdf}

\subsection{Kugelkondensator: Herleitung mit dem Satz von
Gauß}\label{kugelkondensator-herleitung-mit-dem-satz-von-gauuxdf}

\textbf{Elektrisches Feld (Satz von Gauß):}

\[\oint \vec{D} \cdot d\vec{A} = Q\]

\[D(r) \cdot 4\pi r^2 = Q \Rightarrow D(r) = \frac{Q}{4\pi r^2}\]

\[E(r) = \frac{D(r)}{\varepsilon} = \frac{Q}{4\pi \varepsilon r^2}\]

\subsection{Kapazität des
Kugelkondensators}\label{kapazituxe4t-des-kugelkondensators}

\textbf{Spannung zwischen den Kugeln:}

\[U = \int_{R_1}^{R_2} E(r)\,dr = \frac{Q}{4\pi \varepsilon} \int_{R_1}^{R_2} \frac{1}{r^2}\,dr\]

\[U = \frac{Q}{4\pi \varepsilon} \left( \frac{1}{R_1} - \frac{1}{R_2} \right)\]

\textbf{Kapazität:}

\[C = \frac{Q}{U} = 4\pi \varepsilon \frac{R_1 R_2}{R_2 - R_1}\]

\subsection{Zylinderkondensator}\label{zylinderkondensator}

Ein Zylinderkondensator besteht aus zwei koaxialen leitenden Zylindern
mit den Radien \(R_1\) (innen) und \(R_2\) (außen) und der Länge \(l\).

\includegraphics{/tmp/tmp.sFQRuqDRHE/images/image_11.pdf}

\subsection{Zylinderkondensator: Herleitung mit dem Satz von
Gauß}\label{zylinderkondensator-herleitung-mit-dem-satz-von-gauuxdf}

\textbf{Gesucht:} Elektrisches Feld zwischen den Zylindern

\textbf{Ansatz:} Satz von Gauß mit zylindrischer Gauß'scher Fläche
(Radius \(r\), Länge \(l\))

\[\oint \vec{D} \cdot d\vec{A} = Q_{\text{eingeschlossen}}\]

\textbf{Symmetrie:} Das Feld zeigt radial nach außen, konstant auf
Zylinderflächen

\begin{itemize}
\tightlist
\item
  Mantelfläche:
  \(\vec{D} \parallel d\vec{A} \Rightarrow \vec{D} \cdot d\vec{A} = D \cdot dA\)
\item
  Grundflächen:
  \(\vec{D} \perp d\vec{A} \Rightarrow \vec{D} \cdot d\vec{A} = 0\)
\end{itemize}

\[\oint \vec{D} \cdot d\vec{A} = D(r) \cdot 2\pi r l = Q\]

\subsection{Zylinderkondensator: Elektrisches
Feld}\label{zylinderkondensator-elektrisches-feld}

\textbf{Aus dem Satz von Gauß:}

\[D(r) \cdot 2\pi r l = Q\]

\[D(r) = \frac{Q}{2\pi r l}\]

\[E(r) = \frac{D(r)}{\varepsilon} = \frac{Q}{2\pi \varepsilon r l}\]

\textbf{Ergebnis:} Das elektrische Feld nimmt mit \(\frac{1}{r}\) ab.

\subsection{Zylinderkondensator: Spannung und
Kapazität}\label{zylinderkondensator-spannung-und-kapazituxe4t}

\textbf{Spannung zwischen den Zylindern:}

\[U = \int_{R_1}^{R_2} E(r)\,dr = \frac{Q}{2\pi \varepsilon l} \int_{R_1}^{R_2} \frac{1}{r}\,dr\]

\[U = \frac{Q}{2\pi \varepsilon l} \cdot \ln\frac{R_2}{R_1}\]

\textbf{Kapazität:}

\[C = 2 \cdot \pi \cdot \varepsilon_0 \cdot \varepsilon_r \cdot \frac{l}{\ln\frac{R_2}{R_1}}\]

\subsection{Energie im Kondensator}\label{energie-im-kondensator}

Im elektrischen Feld eines Kondensators ist Energie gespeichert, die bei
Entladung wiedergewonnen werden kann.

Während des Aufladevorgangs nimmt die Spannung mit der Ladung
kontinuierlich zu:

\[U(Q) = \frac{Q}{C}\]

Die beim Aufladen gespeicherte Energie berechnet sich zu:

\[W = \int_{0}^{Q} U(q) \cdot dq = \int_{0}^{Q} \frac{q}{C} \cdot dq = \frac{Q^2}{2C}= \frac{1}{2}  Q  U = \frac{1}{2}  C  U^2\]

\[[W] = [U] \cdot [Q] = \text{V} \cdot \text{C} = \text{V} \cdot \text{A} \cdot \text{s} = \text{W} \cdot \text{s} = \text{J}\]

\subsection{Parallelschaltung von
Kondensatoren}\label{parallelschaltung-von-kondensatoren}

Bei der Parallelschaltung von Kondensatoren addieren sich die
Kapazitäten:

\[C_{\text{ges}} = C_1 + C_2 + \dots + C_n = \sum_{i=1}^n C_i\]

Eigenschaften:

\begin{itemize}
\tightlist
\item
  Gleiche Spannung an allen Kondensatoren
\item
  Die Gesamtladung ist die Summe der Einzelladungen
\end{itemize}

\[U_{\text{ges}} = U_1 = U_2 = \dots = U_n\]

\includegraphics{/tmp/tmp.sFQRuqDRHE/images/image_12.pdf}

\subsection{Reihenschaltung von
Kondensatoren}\label{reihenschaltung-von-kondensatoren}

Bei der Reihenschaltung von Kondensatoren addieren sich die Kehrwerte
der Kapazitäten:

\[\frac{1}{C_{\text{ges}}} = \frac{1}{C_1} + \frac{1}{C_2} + \dots + \frac{1}{C_n} = \sum_{i=1}^n \frac{1}{C_i}\]

Eigenschaften:

\begin{itemize}
\tightlist
\item
  Gleiche Ladung auf allen Kondensatoren
  \(Q_1 = Q_2 = \dots = Q_n = Q_{\text{ges}}\)
\item
  Die Gesamtspannung ist die Summe der Einzelspannungen
\end{itemize}

\[U_{\text{ges}} = U_1 + U_2 + \dots + U_n = \sum_{i=1}^n U_i\]

\[C_{\text{ges}} = \frac{Q_{\text{ges}}}{U_{\text{ges}}} = \frac{Q_{\text{ges}}}{U_1 + U_2 + \dots + U_n} =  \frac{Q_{\text{ges}}}{Q_{\text{ges}} \cdot \left(\frac{1}{C_1} + \frac{1}{C_2} + \dots + \frac{1}{C_n}\right)} = \frac{1}{\frac{1}{C_1} + \frac{1}{C_2} + \dots + \frac{1}{C_n}}\]

\subsection{Kondensatoren mit inhomogenen Dielektrika
1}\label{kondensatoren-mit-inhomogenen-dielektrika-1}

Wenn ein Plattenkondensator aus zwei Bereichen mit unterschiedlichen
Dielektrika besteht, berechnet sich die Gesamtkapazität als:

\[C = C_1 + C_2 = \frac{\varepsilon_0}{d} \cdot \left(\varepsilon_{r1} \cdot A_1 + \varepsilon_{r2} \cdot A_2\right)\]

Dies entspricht einer Parallelschaltung von zwei Teilkondensatoren.

\subsection{Kondensatoren mit inhomogenen Dielektrika
2}\label{kondensatoren-mit-inhomogenen-dielektrika-2}

Wenn ein Plattenkondensator aus zwei hintereinander liegenden Schichten
mit unterschiedlichen Dielektrika besteht, berechnet sich die
Gesamtkapazität als:

\[C = \frac{\varepsilon_0 \cdot \varepsilon_{r1} \cdot \varepsilon_{r2} \cdot A}{\varepsilon_{r2} \cdot d_1 + \varepsilon_{r1} \cdot d_2}\]

Dies entspricht einer Reihenschaltung von zwei Teilkondensatoren.

\subsection{Übersicht: Größen im elektrischen
Feld}\label{uxfcbersicht-gruxf6uxdfen-im-elektrischen-feld}

\begin{longtable}[]{@{}
  >{\raggedright\arraybackslash}p{(\columnwidth - 4\tabcolsep) * \real{0.3333}}
  >{\raggedright\arraybackslash}p{(\columnwidth - 4\tabcolsep) * \real{0.3333}}
  >{\raggedright\arraybackslash}p{(\columnwidth - 4\tabcolsep) * \real{0.3333}}@{}}
\toprule\noalign{}
\begin{minipage}[b]{\linewidth}\raggedright
Größe
\end{minipage} & \begin{minipage}[b]{\linewidth}\raggedright
Definition
\end{minipage} & \begin{minipage}[b]{\linewidth}\raggedright
Einheit
\end{minipage} \\
\midrule\noalign{}
\endhead
\bottomrule\noalign{}
\endlastfoot
Elektrische Ladung (\emph{electric charge}) & \(Q\) &
\([Q] = \text{C}\) \\
Spannung (\emph{voltage}) & \(U = \Delta \varphi\) &
\([U] = \text{V}\) \\
Kapazität (\emph{capacitance}) & \(C = \frac{Q}{U}\) &
\([C] = \text{F} = \frac{\text{C}}{\text{V}}\) \\
Elektrische Feldstärke (\emph{electric field {[}strength{]}}) &
\(\vec{E} = \frac{\vec{F}}{Q}\) &
\([\vec{E}] = \frac{\text{V}}{\text{m}}=\frac{\text{N}}{\text{C}}\) \\
Elektrische Flussdichte (\emph{electric flux density}) =
{[}Di-{]}Elektrische Verschiebungsdichte (\emph{electric displacement
field}) & \(\vec{D} = \varepsilon_0 \varepsilon_r \vec{E}\) &
\([\vec{D}] = \frac{\text{C}}{\text{m}^2}\) \\
Elektrische Feldkonstante (\emph{electric constant}) = Permittivität des
Vakuums (\emph{vacuum permittivity}) & \(\varepsilon_0\) &
\([\varepsilon_0] = \frac{\text{C}^2}{\text{N} \cdot \text{m}^2}\) \\
{[}Absolute{]} Permittivität (\emph{{[}absolute{]} permittivity}) =
\st{Dielektrizitätskonstante} & \(\varepsilon\) &
\([\varepsilon] = \frac{\text{C}^2}{\text{N} \cdot \text{m}^2}\) \\
Relative Permittivität (\emph{relative permittivity}) = \st{Relative
Dielektrizitätskonstante} &
\(\varepsilon_r = \frac{\varepsilon}{\varepsilon_0}\) & dimensionslos \\
\end{longtable}

\section{Quiz: Das Elektrische Feld}\label{quiz-das-elektrische-feld}

\subsection{Warum spielt die Gravitation im atomaren Maßstab praktisch
keine
Rolle?}\label{warum-spielt-die-gravitation-im-atomaren-mauxdfstab-praktisch-keine-rolle}

\begin{itemize}
\tightlist
\item
  \begin{enumerate}
  \def\labelenumi{\Alph{enumi})}
  \tightlist
  \item
    Weil Protonen und Elektronen keine Masse besitzen\\
  \end{enumerate}
\item
  \begin{enumerate}
  \def\labelenumi{\Alph{enumi})}
  \setcounter{enumi}{1}
  \tightlist
  \item
    Weil die elektromagnetischen Kräfte viel stärker sind als die
    Gravitationskräfte\\
  \end{enumerate}
\item
  \begin{enumerate}
  \def\labelenumi{\Alph{enumi})}
  \setcounter{enumi}{2}
  \tightlist
  \item
    Weil Gravitation nur zwischen Himmelskörpern wirkt\\
  \end{enumerate}
\item
  \begin{enumerate}
  \def\labelenumi{\Alph{enumi})}
  \setcounter{enumi}{3}
  \tightlist
  \item
    Weil sie durch Quantenmechanik verboten wird
  \end{enumerate}
\end{itemize}

\subsection{Warum sind Atome trotz geladener Bestandteile nach außen
elektrisch
neutral?}\label{warum-sind-atome-trotz-geladener-bestandteile-nach-auuxdfen-elektrisch-neutral}

\begin{itemize}
\tightlist
\item
  \begin{enumerate}
  \def\labelenumi{\Alph{enumi})}
  \tightlist
  \item
    Weil sich Protonen und Neutronen ausgleichen\\
  \end{enumerate}
\item
  \begin{enumerate}
  \def\labelenumi{\Alph{enumi})}
  \setcounter{enumi}{1}
  \tightlist
  \item
    Weil sich die gleiche Anzahl an Protonen (positiv) und Elektronen
    (negativ) kompensiert\\
  \end{enumerate}
\item
  \begin{enumerate}
  \def\labelenumi{\Alph{enumi})}
  \setcounter{enumi}{2}
  \tightlist
  \item
    Weil Elektronen keine Ladung haben\\
  \end{enumerate}
\item
  \begin{enumerate}
  \def\labelenumi{\Alph{enumi})}
  \setcounter{enumi}{3}
  \tightlist
  \item
    Weil neutrale Teilchen überwiegen
  \end{enumerate}
\end{itemize}

\subsection{Was bedeutet, dass elektrische Ladung „quantisiert``
ist?}\label{was-bedeutet-dass-elektrische-ladung-quantisiert-ist}

\begin{itemize}
\tightlist
\item
  \begin{enumerate}
  \def\labelenumi{\Alph{enumi})}
  \tightlist
  \item
    Ladung existiert nur bei Quarks\\
  \end{enumerate}
\item
  \begin{enumerate}
  \def\labelenumi{\Alph{enumi})}
  \setcounter{enumi}{1}
  \tightlist
  \item
    Ladung tritt nur in ganzzahligen Vielfachen einer kleinsten Einheit
    auf\\
  \end{enumerate}
\item
  \begin{enumerate}
  \def\labelenumi{\Alph{enumi})}
  \setcounter{enumi}{2}
  \tightlist
  \item
    Ladung kann beliebige Werte annehmen\\
  \end{enumerate}
\item
  \begin{enumerate}
  \def\labelenumi{\Alph{enumi})}
  \setcounter{enumi}{3}
  \tightlist
  \item
    Ladung hängt vom Beobachter ab
  \end{enumerate}
\end{itemize}

\subsection{Was passiert mit der Kraft zwischen zwei Punktladungen nach
dem Coulomb-Gesetz, wenn der Abstand halbiert
wird?}\label{was-passiert-mit-der-kraft-zwischen-zwei-punktladungen-nach-dem-coulomb-gesetz-wenn-der-abstand-halbiert-wird}

\begin{itemize}
\tightlist
\item
  \begin{enumerate}
  \def\labelenumi{\Alph{enumi})}
  \tightlist
  \item
    Sie halbiert sich\\
  \end{enumerate}
\item
  \begin{enumerate}
  \def\labelenumi{\Alph{enumi})}
  \setcounter{enumi}{1}
  \tightlist
  \item
    Sie vervierfacht sich\\
  \end{enumerate}
\item
  \begin{enumerate}
  \def\labelenumi{\Alph{enumi})}
  \setcounter{enumi}{2}
  \tightlist
  \item
    Sie verdoppelt sich\\
  \end{enumerate}
\item
  \begin{enumerate}
  \def\labelenumi{\Alph{enumi})}
  \setcounter{enumi}{3}
  \tightlist
  \item
    Sie bleibt gleich
  \end{enumerate}
\end{itemize}

\subsection{Was zeigt die Richtung einer elektrischen Feldlinie
an?}\label{was-zeigt-die-richtung-einer-elektrischen-feldlinie-an}

\begin{itemize}
\tightlist
\item
  \begin{enumerate}
  \def\labelenumi{\Alph{enumi})}
  \tightlist
  \item
    Die Bewegung einer negativen Probeladung\\
  \end{enumerate}
\item
  \begin{enumerate}
  \def\labelenumi{\Alph{enumi})}
  \setcounter{enumi}{1}
  \tightlist
  \item
    Die Richtung der Kraft auf eine positive Probeladung\\
  \end{enumerate}
\item
  \begin{enumerate}
  \def\labelenumi{\Alph{enumi})}
  \setcounter{enumi}{2}
  \tightlist
  \item
    Die Richtung der Polarisation\\
  \end{enumerate}
\item
  \begin{enumerate}
  \def\labelenumi{\Alph{enumi})}
  \setcounter{enumi}{3}
  \tightlist
  \item
    Die Richtung minimaler Energie
  \end{enumerate}
\end{itemize}

\subsection{Wann ist der Einsatz des Gaußschen Gesetzes besonders
sinnvoll?}\label{wann-ist-der-einsatz-des-gauuxdfschen-gesetzes-besonders-sinnvoll}

\begin{itemize}
\tightlist
\item
  \begin{enumerate}
  \def\labelenumi{\Alph{enumi})}
  \tightlist
  \item
    Bei beliebigen Ladungsverteilungen\\
  \end{enumerate}
\item
  \begin{enumerate}
  \def\labelenumi{\Alph{enumi})}
  \setcounter{enumi}{1}
  \tightlist
  \item
    Bei jeder einzelnen Punktladung\\
  \end{enumerate}
\item
  \begin{enumerate}
  \def\labelenumi{\Alph{enumi})}
  \setcounter{enumi}{2}
  \tightlist
  \item
    Bei Systemen mit hoher Symmetrie (z. B. Kugel, Zylinder, Ebene)\\
  \end{enumerate}
\item
  \begin{enumerate}
  \def\labelenumi{\Alph{enumi})}
  \setcounter{enumi}{3}
  \tightlist
  \item
    Nur bei negativ geladenen Objekten
  \end{enumerate}
\end{itemize}

\subsection{Warum schwächt ein Dielektrikum das elektrische
Feld?}\label{warum-schwuxe4cht-ein-dielektrikum-das-elektrische-feld}

\begin{itemize}
\tightlist
\item
  \begin{enumerate}
  \def\labelenumi{\Alph{enumi})}
  \tightlist
  \item
    Weil es freie Elektronen enthält\\
  \end{enumerate}
\item
  \begin{enumerate}
  \def\labelenumi{\Alph{enumi})}
  \setcounter{enumi}{1}
  \tightlist
  \item
    Weil es durch Polarisation ein Gegenfeld erzeugt\\
  \end{enumerate}
\item
  \begin{enumerate}
  \def\labelenumi{\Alph{enumi})}
  \setcounter{enumi}{2}
  \tightlist
  \item
    Weil es Ladung vollständig abschirmt\\
  \end{enumerate}
\item
  \begin{enumerate}
  \def\labelenumi{\Alph{enumi})}
  \setcounter{enumi}{3}
  \tightlist
  \item
    Weil es die Permittivität verringert
  \end{enumerate}
\end{itemize}

\subsection{Warum ist die elektrische Arbeit beim Bewegen einer Ladung
im elektrostatischen Feld
wegunabhängig?}\label{warum-ist-die-elektrische-arbeit-beim-bewegen-einer-ladung-im-elektrostatischen-feld-wegunabhuxe4ngig}

\begin{itemize}
\tightlist
\item
  \begin{enumerate}
  \def\labelenumi{\Alph{enumi})}
  \tightlist
  \item
    Weil das Feld nur innerhalb von Leitern existiert\\
  \end{enumerate}
\item
  \begin{enumerate}
  \def\labelenumi{\Alph{enumi})}
  \setcounter{enumi}{1}
  \tightlist
  \item
    Weil Feldlinien immer geschlossen sind\\
  \end{enumerate}
\item
  \begin{enumerate}
  \def\labelenumi{\Alph{enumi})}
  \setcounter{enumi}{2}
  \tightlist
  \item
    Weil es sich um ein Potentialfeld handelt
  \end{enumerate}
\item
  \begin{enumerate}
  \def\labelenumi{\Alph{enumi})}
  \setcounter{enumi}{3}
  \tightlist
  \item
    Weil die Kraft immer konstant ist
  \end{enumerate}
\end{itemize}

\subsection{Was beschreibt das elektrische Potential
physikalisch?}\label{was-beschreibt-das-elektrische-potential-physikalisch}

\begin{itemize}
\tightlist
\item
  \begin{enumerate}
  \def\labelenumi{\Alph{enumi})}
  \tightlist
  \item
    Die Anzahl der Feldlinien\\
  \end{enumerate}
\item
  \begin{enumerate}
  \def\labelenumi{\Alph{enumi})}
  \setcounter{enumi}{1}
  \tightlist
  \item
    Die potenzielle Energie pro Ladung\\
  \end{enumerate}
\item
  \begin{enumerate}
  \def\labelenumi{\Alph{enumi})}
  \setcounter{enumi}{2}
  \tightlist
  \item
    Die Feldstärke unabhängig vom Ort\\
  \end{enumerate}
\item
  \begin{enumerate}
  \def\labelenumi{\Alph{enumi})}
  \setcounter{enumi}{3}
  \tightlist
  \item
    Die Stärke der Polarisation
  \end{enumerate}
\end{itemize}

\section{Gleichstrom}\label{gleichstrom}

\begin{enumerate}
\def\labelenumi{\arabic{enumi}.}
\tightlist
\item
  Stromstärke und Stromdichte
\item
  Widerstand und Ohm'sches Gesetz
\item
  \hyperref[stromkreisberechnungen]{Stromkreisberechnungen}
  (Kirchhoff'sche Regeln)
\item
  Zweipoltheorie
\item
  Arbeit \& Leistung
\end{enumerate}

\subsection{\texorpdfstring{Elektrischer Strom (\emph{electric
current})}{Elektrischer Strom (electric current)}}\label{elektrischer-strom-electric-current}

Strom ist der gerichtete Fluss von elektrischer Ladung

\begin{itemize}
\tightlist
\item
  Stromdichte \(\vec{J} = \rho \cdot \vec{v}\)

  \begin{itemize}
  \tightlist
  \item
    \(\vec{v}\): Geschwindigkeit \emph{positiver} Ladungsträger
  \end{itemize}
\item
  Stromstärke \(I = \int_A \vec{J} \cdot d\vec{A} = \dfrac{dQ}{dt}\)
\item
  \([I] = \text{A} = \dfrac{\text{C}}{\text{s}}\)
\item
  \([\vec{J}] = \dfrac{\text{A}}{\text{m}^2}\)
\end{itemize}

\includegraphics{/tmp/tmp.sFQRuqDRHE/images/image_13.pdf}

\subsection{Stromrichtung \&
Ladungsträger}\label{stromrichtung-ladungstruxe4ger}

\begin{itemize}
\tightlist
\item
  \(\vec{J} = \rho \cdot \vec{v}\) zeigt in die Richtung, in die sich
  \emph{positive} Ladung bewegt -- egal ob die tatsächlichen
  Ladungsträger positiv oder negativ sind!
\item
  Das ist auch die \emph{Zählrichtung} der Stromstärke \(I\)
\end{itemize}

\subsection{Stromleitung in Metallen}\label{stromleitung-in-metallen}

\begin{itemize}
\tightlist
\item
  In Metallen gibt jedes Atom Elektronen ab, die sich frei im Gitter der
  positiv geladenen Atomrümpfe bewegen können („Elektronengas``)
\item
  Die Ladungsdichte der Elektronen ist jederzeit konstant, da eine
  Ansammlung ein elektrisches Feld erzeugen würde, dass durch Abstoßung
  der Elektronen wieder ausgeglichen wird → der Leiter ist überall
  elektrisch neutral
\end{itemize}

\includegraphics{/tmp/tmp.sFQRuqDRHE/images/image_14.pdf}

\subsection{Metalle im elektrischen
Feld}\label{metalle-im-elektrischen-feld}

Klassisches Bild: erfährt das Elektronengas ein elektrisches Feld,
werden die Elektronen beschleunigt, nach kurzer Zeit aber durch Stöße
mit dem Metallgitter wieder abgebremst. Im Mittel ergibt sich dadurch
eine konstante mittleren Geschwindigkeit \(\vec{v}_d\), die
Driftgeschwindigkeit. Sie geht \emph{entgegen} der Feldrichtung
\(\vec{v}_d = \vec{J}/\rho\), \(\rho =- ne\).

\includegraphics{/tmp/tmp.sFQRuqDRHE/images/image_15.pdf}

\subsection{Zahlenbeispiel: Driftgeschwindigkeit im
Kupferdraht}\label{zahlenbeispiel-driftgeschwindigkeit-im-kupferdraht}

Kupfer, \(A=1 \, \text{mm}^2\), \(I=1 \, \text{A}\):

\begin{itemize}
\tightlist
\item
  Dichte freier Elektronen:
  \(n \approx 8{,}5 \cdot 10^{28} \, \frac{1}{\text{m}^3}\)
\item
  Ladungsträgerdichte:
  \(\rho = -n \cdot e \approx -1{,}36 \cdot 10^{10} \, \frac{\text{C}}{\text{m}^3}\)
\item
  Stromdichte:
  \(|\vec{J}| = \frac{I}{A} = \frac{1 \, \text{A}}{1 \cdot 10^{-6} \, \text{m}^2} = 1 \cdot 10^{6} \, \frac{\text{A}}{\text{m}^2}\)
\item
  Driftgeschwindigkeit:
  \(|\vec{v_d}| = \frac{|\vec{J}|}{|\rho|} \approx 7{,}35 \cdot 10^{-5} \, \frac{\text{m}}{\text{s}} \approx 0{,}26 \, \frac{\text{m}}{\text{h}}\)
  🐌
\end{itemize}

\subsection{\texorpdfstring{Elektrische Leitfähigkeit von Metallen
(\emph{electric
conductivity})}{Elektrische Leitfähigkeit von Metallen (electric conductivity)}}\label{elektrische-leitfuxe4higkeit-von-metallen-electric-conductivity}

\begin{itemize}
\tightlist
\item
  Erfährt das Elektronengas ein elektrisches Feld, bewegen sich die
  Elektronen \emph{entgegen} der Feldrichtung
\item
  Für ein gegebenes Material ist die Stromdichte umso höher, je höher
  das elektrische Feld ist
\item
  Der Proportionalitätsfaktor ist die elektrische Leitfähigkeit
  \(\sigma\) des Materials
\end{itemize}

\[\vec{J} = \sigma \cdot \vec{E}\]

= Ohm'sches Gesetz (\emph{Ohm's law})

Achtung: die proportionale Beziehung gilt nur für \emph{lineare Leiter}
(z.B. Metalle bei konstanter Temperatur)

\subsection{Ohm'sches Gesetz im linearen
Leiter}\label{ohmsches-gesetz-im-linearen-leiter}

\[\vec{J} = \sigma \cdot \vec{E}\]

\begin{itemize}
\tightlist
\item
  Stromdichte muss konstant sein
\item
  Elektrisches Feld muss konstant sein
\item
  Potential \(\varphi\) muss linear abfallen

  \begin{itemize}
  \tightlist
  \item
    \(\Phi(l) = \phi(0) - E \cdot l\)
  \item
    \(U = \phi(0) - \phi(l) = E \cdot l\)
  \end{itemize}
\end{itemize}

\[I = J \cdot A = \sigma \cdot E \cdot A = \sigma \cdot \frac{U}{l} \cdot A = \frac{U}{R}\]

\begin{itemize}
\tightlist
\item
  Elektrischer Widerstand \(R = \frac{l}{\sigma \cdot A}\)
  (\emph{electric resistance})
\end{itemize}

\includegraphics{/tmp/tmp.sFQRuqDRHE/images/image_17.pdf}

\subsection{Widerstand und Leitwert}\label{widerstand-und-leitwert}

Der elektrische Widerstand \(R\) ist definiert durch das Ohm'sche
Gesetz:

\[R = \frac{U}{I}\]

Einheit: \([R] = \frac{\text{V}}{\text{A}} = \Omega\) (Ohm)

Der elektrische Leitwert \(G\) ist der Kehrwert des Widerstands:

\[G = \frac{1}{R} = \frac{I}{U}\]

Einheit: \([G] = \frac{\text{A}}{\text{V}} = \text{S}\) (Siemens)

\subsection{Materialeigenschaften
vs.~Bauteilgrößen}\label{materialeigenschaften-vs.-bauteilgruxf6uxdfen}

\textbf{Materialeigenschaften} (intensiv, unabhängig von Geometrie): -
\textbf{Spezifischer Widerstand} \(\rho\): Widerstand pro Längeneinheit
bei Einheitsquerschnitt - \textbf{Leitfähigkeit}
\(\sigma = \frac{1}{\rho}\): Leitfähigkeit des Materials

\textbf{Bauteilgrößen} (extensiv, abhängig von Geometrie): -
\textbf{Widerstand} \(R = \rho \frac{l}{A}\): Widerstand eines konkreten
Leiters - \textbf{Leitwert} \(G = \frac{1}{R} = \sigma \frac{A}{l}\):
Leitwert eines konkreten Leiters

\textbf{Beispiel:} Kupfer hat immer die gleiche Leitfähigkeit
\(\sigma_{\text{Cu}} = 58 \text{ MS/m}\), aber ein dickeres Kabel hat
einen kleineren Widerstand \(R\).

\includegraphics{/tmp/tmp.sFQRuqDRHE/images/image_17.pdf}

\subsection{Übersicht der Größen im linearen
Leiter}\label{uxfcbersicht-der-gruxf6uxdfen-im-linearen-leiter}

\begin{longtable}[]{@{}
  >{\raggedright\arraybackslash}p{(\columnwidth - 6\tabcolsep) * \real{0.2500}}
  >{\raggedright\arraybackslash}p{(\columnwidth - 6\tabcolsep) * \real{0.2500}}
  >{\raggedright\arraybackslash}p{(\columnwidth - 6\tabcolsep) * \real{0.2500}}
  >{\raggedright\arraybackslash}p{(\columnwidth - 6\tabcolsep) * \real{0.2500}}@{}}
\toprule\noalign{}
\begin{minipage}[b]{\linewidth}\raggedright
Größe
\end{minipage} & \begin{minipage}[b]{\linewidth}\raggedright
Definition
\end{minipage} & \begin{minipage}[b]{\linewidth}\raggedright
Einheit
\end{minipage} & \begin{minipage}[b]{\linewidth}\raggedright
Name
\end{minipage} \\
\midrule\noalign{}
\endhead
\bottomrule\noalign{}
\endlastfoot
Spannung (\emph{voltage}) & \(U = \Delta \varphi\) & \([U] = \text{V}\)
& Volt \\
Stromstärke (\emph{current}) & \(I = \frac{\Delta Q}{\Delta t}\) &
\([I] = \text{A}\) & Ampere \\
Widerstand (\emph{resistance}) & \(R = \frac{U}{I}\) & \([R] = \Omega\)
& Ohm \\
Leitwert (\emph{conductance}) & \(G = \frac{1}{R}\) &
\([G] = \text{S} = \frac{1}{\Omega}\) & Siemens \\
spezifischer Widerstand (\emph{resistivity}) & \(\rho = R \frac{A}{l}\)
& \([\rho] = \Omega \cdot \text{m}\) & Ohm-Meter \\
Leitfähigkeit (\emph{conductivity}) & \(\sigma = \frac{1}{\rho}\) &
\([\sigma] = \text{S/m}\) & Siemens pro Meter \\
\end{longtable}

\subsection{Temperaturabhängigkeit des
Widerstands}\label{temperaturabhuxe4ngigkeit-des-widerstands}

Bei den meisten Materialien ändert sich der Widerstand mit der
Temperatur.

Kleinsignalverhalten (lineare Näherung):

\[R(T) = R(T_0) \cdot [1 + \alpha \cdot (T - T_0)]\]

Dabei ist:

\begin{itemize}
\tightlist
\item
  \(\alpha\) der Temperaturkoeffizient des Widerstands (Einheit:
  \([\alpha] = \frac{1}{\text{K}}\))
\item
  \(T_0\) die Bezugstemperatur (üblicherweise 20°C oder 0°C)
\item
  \(T\) die aktuelle Temperatur
\end{itemize}

\subsection{Elektrische Leitfähigkeit verschiedener
Materialien}\label{elektrische-leitfuxe4higkeit-verschiedener-materialien}

Bei Leitern nimmt der Widerstand mit steigender Temperatur zu (positiver
Temperaturkoeffizient α \textgreater{} 0).

Typische Werte für einige Leitermaterialien bei 20°C: \textbar{}
Leitermaterial \textbar{} Spez. Widerstand \(\rho\) (µΩ·m) \textbar{}
Leitfähigkeit \(\sigma\) (MS/m) \textbar{} Temperaturkoeffizient
\(\alpha\) (1/K) \textbar{}
\textbar----------------\textbar--------------------------------------\textbar-------------------------\textbar--------------------------------------\textbar{}
\textbar{} Silber \textbar{} 0,016 \textbar{} 63 \textbar{} 3,8 · 10−3
\textbar{} \textbar{} Kupfer \textbar{} 0,017 \textbar{} 58 \textbar{}
3,9 · 10−3 \textbar{} \textbar{} Aluminium \textbar{} 0,027 \textbar{}
38 \textbar{} 4,3 · 10−3 \textbar{} \textbar{} Messing \textbar{} 0,062
\textbar{} 16 \textbar{} 2,0 · 10−3 \textbar{}

\(R(T) = R(T_0) \cdot [1 + \alpha \cdot (T - T_0)]\)

\subsection{Metalle als
Temperatursensoren}\label{metalle-als-temperatursensoren}

Die Temperaturabhängigkeit des Widerstands macht Metalle zu präzisen
Temperatursensoren.

\textbf{Platin-Widerstandsthermometer (Pt100):} - Pt100:
\(R(0°\text{C}) = 100 \, \Omega\) - Temperaturkoeffizient:
\(\alpha_{\text{Pt}} = 3{,}85 \cdot 10^{-3} \, \text{K}^{-1}\)

\textbf{Vorteile von Platin-Sensoren:} - Hohe Langzeitstabilität -
Breiter Messbereich (-200°C bis +850°C) - Gute Linearität - Chemische
Beständigkeit

\includegraphics{/tmp/tmp.sFQRuqDRHE/images/a835be2372c6400cc73d95b00dd3007f.jpg}

\section{Stromkreisberechnungen}\label{stromkreisberechnungen}

\begin{enumerate}
\def\labelenumi{\arabic{enumi}.}
\tightlist
\item
  Die Kirchhoff'schen Gesetze
\item
  Zweipoltheorie
\item
  Arbeit und Leistung im Gleichstromkreis
\end{enumerate}

\subsection{Knotenpunktregel (1. Kirchhoff'sches
Gesetz)}\label{knotenpunktregel-1.-kirchhoffsches-gesetz}

In einem Knotenpunkt kann weder Ladung gespeichert noch erzeugt werden.
Die Summe aller zufließenden Ströme ist gleich der Summe aller
abfließenden Ströme:

\[\sum_{k} I_{k} = 0\]

\includegraphics{/tmp/tmp.sFQRuqDRHE/images/image_19.pdf}

\subsection{Mathematische Analogie}\label{mathematische-analogie}

In einer stationären (zeitlich unveränderlichen) Stromverteilung ist die
elektrische Ladung überall konstant. Der gesamte Strom durch jede
geschlossene Oberfläche ist Null:

\[\oint_{A} \vec{J} \cdot d\vec{A} = - \frac{dQ_{\text{innen}}}{dt} = 0\]

Vgl. Satz von Gauß in Abwesenheit von eingeschlossener Ladung:

\[\oint_{A} \vec{D} \cdot d\vec{A} = Q_{\text{innen}} = 0\]

(NB: die obige Gleichung folgt \emph{nicht} aus der unteren -- die
mathematische Analogie gilt, da sowohl das elektrostatische Feld als
auch die stationäre Stromdichte \textbf{quellenfreie Vektorfelder}
sind.)

\subsection{Maschenregel (2. Kirchhoff'sches
Gesetz)}\label{maschenregel-2.-kirchhoffsches-gesetz}

Die Summe aller in einer Masche auftretenden Spannungen ist Null:

\[\sum_{k} U_{k} = 0\]

\includegraphics{/tmp/tmp.sFQRuqDRHE/images/image_20.pdf}

\subsection{Reihenschaltung von
Widerständen}\label{reihenschaltung-von-widerstuxe4nden}

\[R_\text{ges} = R_1 + R_2 + R_3 + \dots + R_n = \sum_{i=1}^{n} R_i\]

\begin{itemize}
\tightlist
\item
  Gleicher Strom durch alle Widerstände: \(I = I_1 = I_2 = \dots = I_n\)
\item
  Die Gesamtspannung ist die Summe der Einzelspannungen:
  \(U_\text{ges} = U_1 + U_2 + \dots + U_n\)
\item
  Gesamtwiderstand ist größer als der größte Einzelwiderstand
\end{itemize}

\includegraphics{/tmp/tmp.sFQRuqDRHE/images/image_21.pdf}

\subsection{Spannungsteiler}\label{spannungsteiler}

Bei einer Reihenschaltung teilt sich die Gesamtspannung im Verhältnis
der Widerstände auf:

\[I = \frac{U}{R_\text{ges}} = \frac{U_1}{R_1} = \frac{U_2}{R_2}\]

\includegraphics{/tmp/tmp.sFQRuqDRHE/images/image_22.pdf}

\subsection{Parallelschaltung von
Widerständen}\label{parallelschaltung-von-widerstuxe4nden}

Bei einer Parallelschaltung von Widerständen addieren sich die Leitwerte
zum Gesamtleitwert:

\[\frac{1}{R_\text{ges}} = \frac{1}{R_1} + \frac{1}{R_2} + \dots + \frac{1}{R_n} = \sum_{i=1}^{n} \frac{1}{R_i}\]

Oder mit Leitwerten:

\[G_\text{ges} = G_1 + G_2 + \dots + G_n = \sum_{i=1}^{n} G_i\]

\begin{itemize}
\tightlist
\item
  Der Gesamtstrom ist die Summe der Einzelströme:
  \(I_\text{ges} = I_1 + I_2 + \dots + I_n\)
\item
  Gleiche Spannung an allen Widerständen:
  \(U = U_1 = U_2 = \dots = U_n\)
\item
  Der Gesamtwiderstand ist kleiner als der kleinste Einzelwiderstand
\end{itemize}

\includegraphics{/tmp/tmp.sFQRuqDRHE/images/image_23.pdf}

\subsection{Parallelschaltung:
Herleitung}\label{parallelschaltung-herleitung}

Wegen der Knotenregel gilt:

\[I_\text{ges} = I_1 + I_2 + \dots + I_n = \frac{U}{R_1} + \frac{U}{R_2} + \dots + \frac{U}{R_n}\]

Außerdem per Definition:

\[I_\text{ges} = \frac{U}{R_\text{ges}}\]

Es folgt:

\[\Rightarrow\frac{U}{R_\text{ges}} = U \cdot \left(\frac{1}{R_1} + \frac{1}{R_2} + \dots + \frac{1}{R_n}\right)\]

\[\Rightarrow\frac{1}{R_\text{ges}} = \frac{1}{R_1} + \frac{1}{R_2} + \dots + \frac{1}{R_n}\]

\subsection{Stromteilerregel}\label{stromteilerregel}

Bei einer Parallelschaltung teilt sich der Gesamtstrom im umgekehrten
Verhältnis der Widerstände bzw. im direkten Verhältnis der Leitwerte
auf:

\[\frac{I}{G_\text{ges}} = \frac{I_1}{G_1} = \frac{I_2}{G_2}= \dots = \frac{I_n}{G_n}\]

\includegraphics{/tmp/tmp.sFQRuqDRHE/images/image_24.pdf}

\section{Zweipoltheorie}\label{zweipoltheorie}

Ein Zweipol (\emph{two-pole}) oder Eintor (\emph{one-port}) ist ein
elektrisches Bauteil mit zwei zugänglichen Anschlüssen

\includegraphics{/tmp/tmp.sFQRuqDRHE/images/image_25.pdf}

\textbf{Gliederung}

\begin{enumerate}
\def\labelenumi{\arabic{enumi}.}
\tightlist
\item
  Passive lineare Zweipole
\item
  Aktive lineare Zweipole

  \begin{enumerate}
  \def\labelenumii{\arabic{enumii}.}
  \tightlist
  \item
    Ideale Spannungsquelle
  \item
    Ideale Stromquelle
  \item
    Reale Spannungsquelle
  \item
    Reale Stromquelle
  \item
    Äquivalenz von realer Spannungs- und Stromquelle
  \end{enumerate}
\end{enumerate}

\subsection{Passive lineare Zweipole}\label{passive-lineare-zweipole}

\begin{itemize}
\tightlist
\item
  Passiv: Zweipol gibt keine Energie ab
\item
  Linear: Strom-Spannungs-Kennlinie ist eine Gerade
\end{itemize}

Passive lineare Zweipole können zu einem Ersatzwiderstand
zusammengefasst werden

\[U = R \cdot I\]

\subsection{Ideale Spannungsquelle}\label{ideale-spannungsquelle}

Eine ideale Spannungsquelle liefert eine konstante Spannung \(U_0\)
unabhängig von der Belastung.

Eigenschaften:

\begin{itemize}
\tightlist
\item
  Konstante Klemmenspannung \(U = U_0\)
\item
  Innenwiderstand \(R_i = 0\)
\item
  Beliebiger Strom \(I\) möglich
\end{itemize}

\includegraphics{/tmp/tmp.sFQRuqDRHE/images/image_26.pdf}

\subsection{Ideale Stromquelle}\label{ideale-stromquelle}

Eine ideale Stromquelle liefert einen konstanten Strom \(I_0\)
unabhängig von der Belastung.

Eigenschaften:

\begin{itemize}
\tightlist
\item
  Konstanter Strom \(I = I_0\)
\item
  Innenwiderstand \(R_i = \infty\)
\item
  Beliebige Spannung \(U\) möglich
\end{itemize}

\includegraphics{/tmp/tmp.sFQRuqDRHE/images/image_27.pdf}

\subsection{Reale Spannungsquelle}\label{reale-spannungsquelle}

Eine reale Spannungsquelle kann als Reihenschaltung einer idealen
Spannungsquelle \(U_0\) mit einem Innenwiderstand \(R_i\) dargestellt
werden.

Eigenschaften:

\begin{itemize}
\tightlist
\item
  Klemmenspannung nimmt mit zunehmendem Strom ab:
  \(U = U_0 - R_i \cdot I\)
\item
  Bei Leerlauf: \(U = U_0\) (maximale Spannung)
\item
  Bei Kurzschluss: \(I = \frac{U_0}{R_i}\) (maximaler Strom)
\end{itemize}

\includegraphics{/tmp/tmp.sFQRuqDRHE/images/image_30.pdf}

\subsection{Reale Stromquelle}\label{reale-stromquelle}

Eine reale Stromquelle kann als Parallelschaltung einer idealen
Stromquelle \(I_0\) mit einem Innenwiderstand \(R_i\) dargestellt
werden.

Eigenschaften:

\begin{itemize}
\tightlist
\item
  Strom nimmt mit zunehmender Spannung ab: \(I = I_0 - \frac{U}{R_i}\)
\item
  Bei Leerlauf: \(U = I_0 \cdot R_i\) (maximale Spannung)
\item
  Bei Kurzschluss: \(I = I_0\) (maximaler Strom)
\end{itemize}

\includegraphics{/tmp/tmp.sFQRuqDRHE/images/image_31.pdf}

\subsection{Äquivalenz von realer Spannungs- und
Stromquelle}\label{uxe4quivalenz-von-realer-spannungs--und-stromquelle}

Die reale Spannungsquelle und reale Stromquelle sind äquivalent, wenn
folgende Beziehungen gelten:

\[U_0 = I_0 \cdot R_i\]

\[I_0 = \frac{U_0}{R_i}\]

\textbf{Umrechnung:} - Spannungsquelle → Stromquelle:
\(I_0 = \frac{U_0}{R_i}\) - Stromquelle → Spannungsquelle:
\(U_0 = I_0 \cdot R_i\)

Beide Darstellungen beschreiben dieselbe I-U-Kennlinie:
\(U_\text{kl} = U_0 - R_i \cdot I\)

\includegraphics{/tmp/tmp.sFQRuqDRHE/images/image_30.pdf}

\includegraphics{/tmp/tmp.sFQRuqDRHE/images/image_31.pdf}

\subsection{Reihenschaltung von aktiven, linearen
Zweipolen}\label{reihenschaltung-von-aktiven-linearen-zweipolen}

Bei der Reihenschaltung von realen Spannungsquellen addieren sich die
Leerlaufspannungen und die Innenwiderstände:

\[U_{0,\text{ges}} = U_{0,1} + U_{0,2} + \dots + U_{0,n}\]

\[R_{i,\text{ges}} = R_{i,1} + R_{i,2} + \dots + R_{i,n}\]

\textbf{Anwendung:} Batteriepacks in Taschenlampen, Elektroautos
\textbf{Vorteil:} Höhere Gesamtspannung \textbf{Nachteil:} Höherer
Innenwiderstand, bei Ausfall einer Quelle fällt das gesamte System aus

\subsection{Parallelschaltung von aktiven, linearen
Zweipolen}\label{parallelschaltung-von-aktiven-linearen-zweipolen}

Bei der Parallelschaltung von realen Spannungsquellen mit gleicher
Leerlaufspannung \(U_0\) addieren sich die Leitwerte der
Innenwiderstände:

\[\frac{1}{R_{i,\text{ges}}} = \frac{1}{R_{i,1}} + \frac{1}{R_{i,2}} + \dots + \frac{1}{R_{i,n}}\]

Die gemeinsame Leerlaufspannung bleibt \(U_0\).

\textbf{Anwendung:} Notstromversorgung, Batteriepacks für höhere Ströme
\textbf{Vorteil:} Geringerer Innenwiderstand, höhere verfügbare Ströme
\textbf{Nachteil:} Nur bei gleichen Spannungen sinnvoll,
Ausgleichsströme bei unterschiedlichen Quellen

\section{Arbeit und Leistung in
Gleichstromkreisen}\label{arbeit-und-leistung-in-gleichstromkreisen}

\begin{enumerate}
\def\labelenumi{\arabic{enumi}.}
\tightlist
\item
  Elektrische Arbeit
\item
  Elektrische Leistung
\end{enumerate}

\subsection{Elektrische Arbeit
(Energie)}\label{elektrische-arbeit-energie}

Die elektrische Arbeit ist definiert als das Produkt aus Spannung, Strom
und Zeit:

\[W = U \cdot I \cdot t = I^2 \cdot R \cdot t = \frac{U^2}{R} \cdot t\]

Einheit:
\([W] = \text{V} \cdot \text{A} \cdot \text{s} =\text{W} \cdot \text{s} = \text{J}\)
(Joule)

\subsection{Elektrische Leistung}\label{elektrische-leistung}

Die elektrische Leistung ist definiert als elektrische Arbeit pro
Zeiteinheit:

\[P = \frac{dW}{dt} = U \cdot I = I^2 \cdot R = \frac{U^2}{R}\]

Einheit: \([P] = \text{W}\) (Watt)

\section{Leistungsanpassung}\label{leistungsanpassung}

Die Leistungsanpassung beschäftigt sich mit der Frage, bei welchem
Verbraucherwiderstand \(R\) die maximale Leistung aus einer Quelle
entnommen werden kann.

Für eine reale Spannungsquelle mit der Leerlaufspannung \(U_0\) und dem
Innenwiderstand \(R_i\) beträgt die Leistung am Verbraucher:

\[P = R \cdot I^2 = U_0^2 \cdot \frac{R}{(R_i + R)^2}\]

Diese Leistung wird maximal, wenn der Verbraucherwiderstand gleich dem
Innenwiderstand der Quelle ist:

\[R = R_i\]

Die maximale Leistung beträgt dann:

\[P_\text{max} = \frac{U_0^2}{4 \cdot R_i}\]

\subsection{Anpassungsverhältnis und
Wirkungsgrad}\label{anpassungsverhuxe4ltnis-und-wirkungsgrad}

Das Anpassungsverhältnis \(\alpha\) ist definiert als:

\[\alpha = \frac{R}{R_i}\]

Der Wirkungsgrad \(\eta\) gibt das Verhältnis der am Verbraucher
umgesetzten Leistung zur Gesamtleistung der Quelle an:

\[\eta = \frac{P}{P_0} = \frac{R}{R_i + R} = \frac{\alpha}{1 + \alpha}\]

Bei optimaler Leistungsanpassung (\(\alpha = 1\)) beträgt der
Wirkungsgrad nur \(\eta = 0,5\) (50\%).

\subsection{Betriebszustände einer aktiven
Quelle}\label{betriebszustuxe4nde-einer-aktiven-quelle}

\begin{longtable}[]{@{}
  >{\raggedright\arraybackslash}p{(\columnwidth - 8\tabcolsep) * \real{0.1156}}
  >{\raggedright\arraybackslash}p{(\columnwidth - 8\tabcolsep) * \real{0.1497}}
  >{\raggedright\arraybackslash}p{(\columnwidth - 8\tabcolsep) * \real{0.3673}}
  >{\raggedright\arraybackslash}p{(\columnwidth - 8\tabcolsep) * \real{0.2245}}
  >{\raggedright\arraybackslash}p{(\columnwidth - 8\tabcolsep) * \real{0.1429}}@{}}
\toprule\noalign{}
\begin{minipage}[b]{\linewidth}\raggedright
\end{minipage} & \begin{minipage}[b]{\linewidth}\raggedright
Last
\end{minipage} & \begin{minipage}[b]{\linewidth}\raggedright
Leistung Quelle \(P_0\)
\end{minipage} & \begin{minipage}[b]{\linewidth}\raggedright
Leistung Last \(P\)
\end{minipage} & \begin{minipage}[b]{\linewidth}\raggedright
Wirkungsgrad \(\eta\)
\end{minipage} \\
\midrule\noalign{}
\endhead
\bottomrule\noalign{}
\endlastfoot
Kurzschluß & \(R = 0\) & \(P_0 = \frac{U_0^2}{R_i}\) & \(P = 0\) &
\(\eta = 0\) \\
Unteranpassung & \(R < R_i\) &
\(P_0 = \frac{U_0^2}{R_i} \cdot \frac{R}{R+R_i}\) &
\(0 < P < P_\text{max}\) & \(0 < \eta < 0{,}5\) \\
Anpassung & \(R = R_i\) & \(P_0 = \frac{U_0^2}{2R_i}\) &
\(P = \frac{U_0^2}{4R_i}\) & \(\eta = 0{,}5\) \\
Überanpassung & \(R > R_i\) &
\(P_0 = \frac{U_0^2}{R_i} \cdot \frac{R}{R+R_i}\) &
\(0 < P < P_\text{max}\) & \(0{,}5 < \eta < 1\) \\
Leerlauf & \(R \to \infty\) & \(P_0 = 0\) & \(P = 0\) & \(\eta = 1\) \\
\end{longtable}

\end{document}
